\section{Introduction}
Modeling and parameterizing the stably stratified atmospheric boundary layer (SBL) remains one of the most persistent challenges in geophysical fluid dynamics, numerical weather prediction (NWP), and climate modeling \cite{gabls3_2013, sheba2002}. Under strongly stable conditions, turbulent mixing is suppressed by buoyancy forces, leading to a decoupling of the surface from the overlying atmosphere, the formation of nocturnal Low-Level Jets (LLJs), and intense surface radiative cooling \cite{cases99, grachev2007}. Standard Monin--Obukhov Similarity Theory (MOST) closures—originally derived under constant-flux surface layer assumptions—regularly break down in these highly stratified, non-stationary regimes, causing single-column models (SCMs) and Earth System Models (ESMs) to exhibit severe temperature and wind velocity biases \cite{cases99, gabls3_2013, businger1971}.

To resolve these biases, developers conventionally calibrate empirical stability parameters (such as the Businger-Dyer coefficients $\beta_m, \beta_h$) when simulated profiles diverge from observations \cite{gabls3_2013, businger1971, STABILITY}. However, vertical profiles of the gradient Richardson number ($Ri_g$) reconstructed from meteorological towers and remote sensing campaigns frequently exhibit sharp inflection "knees" and apparent coordinate folds \cite{cases99, BelowFastSlow}. In this work, we present a unified mathematical framework grounded in Geometric Singular Perturbation Theory (GSPT) to demonstrate that these inflections are often coordinate-mapping features ("Fold Illusions") generated by the non-linear projection of similarity coordinates under active vertical flux divergence, rather than genuine dynamical transitions of the boundary layer \cite{fenichel1979, CleanRestatement}.

Dynamically, SBL turbulence collapse can be modeled as a multiscale fast-slow dynamical system, where the microscale Turbulent Kinetic Energy (TKE) relaxing rapidly onto a critical manifold $\mathcal{C}_0$ governed by the slower evolution of mean vertical wind shear $S$ \cite{BelowFastSlow, CleanRestatement}. The scale separation between these turbulent and geostrophic subsystems is parameterized by a small dimensionless ratio $0 < \epsilon \ll 1$ \cite{BelowFastSlow}:
\begin{equation}
    \epsilon \frac{dE}{dt} = f(E, S) = l \sqrt{E+\delta} \left( S^2 - \phi N^2 - c_b N^2 \frac{E}{E+\alpha} \right) - \frac{E^{3/2}}{l}, \label{eq:fast_tke_intro}
\end{equation}
\begin{equation}
    \frac{dS}{dt} = g(E, S) = G - c_1 E S - c_2 S, \label{eq:slow_shear_intro}
\end{equation}
where $E$ represents the grid-resolved TKE, $l$ is the turbulent mixing length, $N^2$ is the Brunt-V\"ais\"al\"a buoyancy frequency, and $G$ represents geostrophic wind forcing \cite{BelowFastSlow, CoreMathRevisions}.

A long-standing computational issue in these multiscale models is the existence of a mathematical singularity in the production Jacobian ($\partial f / \partial E$) as TKE approaches zero \cite{CoreMathRevisions}. Standard closures closed with a velocity scale of $\sqrt{E}$ suffer from infinite derivative growth ($\lim_{E \to 0^+} \frac{\partial}{\partial E}(l\sqrt{E}S^2) = +\infty$), triggering severe numerical instabilities near the laminar state \cite{CoreMathRevisions}. To restore mathematical regularity and ensure a globally $C^1$-continuous fast vector field on the physically admissible domain $\mathcal{D} = \{ (E, S) \in \mathbb{R}^2 : E \ge 0 \}$, we embed a velocity-scale regularization floor $\delta = 10^{-6}\text{ m}^2\text{s}^{-2}$ directly into the shear production and buoyant loss terms of the TKE budget \cite{BelowFastSlow, CoreMathRevisions}.

By formulating SBL transitions as slow-passage bifurcations across a regularized critical manifold, we establish a mathematically consistent, singularity-free boundary-layer architecture. This paper is organized as follows: Section 2 derives the continuous differential geometry of SBL coordinate curvature under Nieuwstadt local scaling; Section 3 analyzes the topological bifurcations and hysteresis of the non-monotonic buoyancy closure; Section 4 provides mathematically rigorous proofs for discrete grid spacing and Fenichel manifold slaving convergence; Section 5 outlines the operational implications for SCM and NWP error partitioning; and Section 6 synthesizes our findings and details future research avenues.